\subsection*{An exact obstruction to the intermediate free-product presentation}

\statement{B.67}{Refuted}{The proposed intermediate braid injection is false}
In B.65's cyclic source, delete moving label $5$ from the
seven-point rigid target.  The resulting map to $PB_5$
kills $f_1$ and factors through
\[
 Q=\langle a,t,c,f,e\mid[a,t]=[c,f]=1\rangle
       \cong\mathbb Z^2*\mathbb Z^2*\mathbb Z,
                                                               \tag{B.67.1}
\]
where
\[
 a=f_0,\quad c=f_2,\quad f=f_5,\quad e=f_4,
 \qquad t=f_5^{-1}f_4^{-1}f_0^{-1}f_2^{-1}f_3^{-1}.
                                                               \tag{B.67.2}
\]
The induced map $Q\to PB_5$ is not injective.
An explicit nonidentity kernel element is
\[
       w=[(ae)^{-1}f(ae),c].
                                                               \tag{B.67.3}
\]
It remains nonidentity in $\widehat Q$ and belongs to
the kernel of $\widehat Q\to\widehat{PB}_5$.
The corresponding source word is, however, nonidentity
under the original six-strand map $\theta_5$.
Thus (B.67.3) refutes the proposed intermediate image
presentation; it does not settle the full cyclic
rigid-kernel comparison.

\begin{proof}
All maps use B.61's actual common geometric marking
and its certified eleven-strand words.
After deleting moving label $5$, the five finite labels
are $(4,11,1,10,6)$.
The direct faithful Artin evaluation recorded in
\nolinkurl{next_joint_rigid_symplectic/reconstruct_Q_birman.py}
checks that $f_1$, $[f_0,t]$, and $[f_2,f_5]$ have
identity images in $PB_5$.
In the free group after filling $f_1$, the five elements
$(a,t,c,f,e)$ form a Nielsen basis, since the inverse
formula is
\[
                 f_3=t^{-1}f^{-1}e^{-1}a^{-1}c^{-1}.
                                                               \tag{B.67.4}
\]
This proves the factorization through (B.67.1).
The meridian $t$ is the infinity meridian after filling
$f_1$; no substitution for the original six-generator
infinity meridian is made.

For a structural derivation of the additional relation,
move label $10$ to the final strand.  The finite order
is then $(4,11,1,6,10)$, and the last-strand Birman
kernel is a free group on four generators $y_1,\ldots,y_4$.
Use the free basis
\[
 C=y_3^{-1}y_2^{-1}y_1y_2y_3,\quad
 T=y_3^{-1}y_2y_3,\quad R=y_3,\quad S=y_4.
                                                               \tag{B.67.5}
\]
The inverse change is
$y_1=RTC T^{-1}R^{-1}$, $y_2=RTR^{-1}$,
$y_3=R$, $y_4=S$.
The exact action certificate gives
\begin{align*}
 a(C)&=(RT)^{-1}C(RT),& a(S)&=(RT)^{-1}S(RT),\\
 a(T)&=T,&a(R)&=R,\\
 f(C)&=C,&f(T)&=STS^{-1},\\
 f(R)&=[S,T]R[S,T]^{-1},&
 f(S)&=STS T^{-1}S^{-1},\\
 e(C)&=S^{-1}CS.&&
                                                               \tag{B.67.6}
\end{align*}
All these are actual conjugation actions in the Birman
fiber, not actions on an associated graded quotient.
The same certificate records the remaining images for
$e$ and checks the action on conjugates of both $c,t$.

Put $V=T^{-1}RT$ and $\beta=ae$.
The formulas for $f$ show that it fixes $C,V,ST$.
With the composition convention
$\alpha_{ae}=\alpha_a\circ\alpha_e$, the two required
formulas for $a$ and the formula for $e(C)$ give
\[
 \beta(C)=pCp^{-1},\qquad
 p=(RT)^{-1}S^{-1}=V^{-1}(ST)^{-1}.
                                                               \tag{B.67.7}
\]
Therefore $f$ fixes $\beta(C)$ and
$\beta^{-1}f\beta$ fixes $C$.
Under the genuine free-fiber homomorphism the source
element $c$ maps to $C^{-1}$, which is fixed as well.
This proves that the actual five-strand image of
(B.67.3) is the identity.

For clarity about conventions, the program that reads
the last Artin conjugator first produces an
\emph{antihomomorphism} on the Birman kernel.
It checks explicitly on two noncommuting canonical
fiber generators that
$\operatorname{push}(hk)=\operatorname{push}(k)
\operatorname{push}(h)$.
Inverting every output supplies the homomorphism and
sends $c$ to $C^{-1}$ rather than $C$.
Conjugation automorphisms commute with inversion, so
this correction preserves the fixed-word argument.
The subsequent direct evaluation of the complete
braid word provides an independent check of the
identity, without relying only on this convention.

The word is nonidentity in the proposed group $Q$.
In the free product $A*B*E$, where
$A=\langle a,t\rangle$, $B=\langle c,f\rangle$,
$E=\langle e\rangle$, its expansion is
\[
 e^{-1}a^{-1}fae\,c\,e^{-1}a^{-1}f^{-1}ae\,c^{-1}.
                                                               \tag{B.67.8}
\]
Its twelve nonidentity syllables lie successively in
$E,A,B,A,E,B,E,A,B,A,E,B$.
This is a nonempty reduced free-product normal form.
There is also a small finite witness: send $a,c,f$
to $(1\ 2)$, send $e$ to $(2\ 3)$, and send $t$
to the identity in $S_3$.
Both relations in (B.67.1) hold, while (B.67.8)
maps to a nontrivial three-cycle.
This proves nonidentity in $\widehat Q$ as well.

The complete fresh calculation is
\nolinkurl{next_joint_rigid_symplectic/refute_Q_injection.py},
with its actual standard output and
\nolinkurl{Q_noninjection_certificate.json}.
With signed index $j$ denoting $f_{j-1}$, the source word is
\[
              (-5,-1,6,1,5,3,-5,-1,-6,1,5,-3).
                                                               \tag{B.67.9}
\]
The reduced five-strand substitution has length $98$;
all five Artin images are the respective single
generators.  The reduced six-strand substitution has
length $166$, and its Artin images have lengths
\[
                   (15705,9489,9119,6417,30465,15129).
                                                               \tag{B.67.10}
\]
The certificate stores the entire first nonidentity
image, as well as every source and braid word used.
The source is a commutator, so its exponent sum is zero.
B.65's fixed-pair linking argument removes the central
full-twist ambiguity.  Consequently (B.67.9) is
nonidentity in the full seven-point sphere target.

In particular the stabilizer of $C$ in the actual base
free group $L=F(a,f,e)$ contains both $f$ and
$(ae)^{-1}f(ae)$; these generate a nonabelian free
subgroup.  Thus the earlier proposed $C$-orbit family
indexed by $L/\langle f\rangle$ already repeats actual
elements, and its proposed free independence is false.
Adding (B.67.3) gives another intermediate quotient,
but a complete presentation of the image and the
full $\theta_5$ comparison are not asserted.
\end{proof}
